\begin{lemma}
Let $\alpha$ satisfy $0<\alpha\le1$, let $\delta>0$ and let $q$ satisfy
\[
  1\le q<\frac{2}{2-\alpha}.
\]
Then the function $z\mapsto|z|^{(\alpha-2)q}$ is integrable on the disk
$\{z:|z|<\delta\}$ for plane Lebesgue measure; equivalently,
$|z|^{\alpha-2}$ belongs to $L^{q}$ of that disk.

This is the assertion the paper makes in one line at paper.tex:534, ``The
function $|z|^{\alpha-2}$ is integrable over the disk $0\le|z|<\delta$ with
every exponent $q$ such that $1\le q<2/(2-\alpha)$'', and it is what turns the
estimate (52) at $m=k+1$ into $u\in W^{k+2,q}_{\mathrm{loc}}(\Omega)$.  The
range of $q$ is exactly (10).
\end{lemma}
\begin{proof}
Write $s=(\alpha-2)q$ and $d\sigma$ for plane Lebesgue measure.

\emph{Step 1: the exponent.}  Since $\alpha\le1$ we have $2-\alpha\ge1>0$, so
the inequality $q<2/(2-\alpha)$ may be multiplied by $2-\alpha>0$ to give
$(2-\alpha)q<2$, that is
\begin{equation}\tag{1}
  s=(\alpha-2)q=-(2-\alpha)q>-2 .
\end{equation}
Equivalently $s+2>0$.  Note also $s<0$, since $\alpha-2<0$ and $q\ge1>0$.

\emph{Step 2: measurability.}  The function $z\mapsto|z|^{s}$ is continuous on
$\mathbb{C}-\{0\}$, being a composition of the continuous $z\mapsto|z|$, which
is positive there, with the continuous $t\mapsto t^{s}$ on $(0,\infty)$; and
$\{0\}$ is a Lebesgue null set.  Hence the function is measurable, and it is
nonnegative, so integrability on the disk is equivalent to the finiteness of
its integral there.

\emph{Step 3: the polar computation.}  Passing to polar coordinates about the
origin, $d\sigma=r\,dr\,d\theta$,
\[
  \int_{|z|<\delta}|z|^{s}\,d\sigma
  =\int_0^{2\pi}\!\!\int_0^{\delta}r^{s}\cdot r\,dr\,d\theta
  =2\pi\int_0^{\delta}r^{s+1}\,dr .
\]
By (1) we have $s+1>-1$, so the one-dimensional integral converges:
\[
  \int_0^{\delta}r^{s+1}\,dr=\frac{\delta^{s+2}}{s+2}<\infty ,
\]
the denominator being positive by (1).  Therefore
\[
  \int_{|z|<\delta}|z|^{(\alpha-2)q}\,d\sigma
  =\frac{2\pi\,\delta^{(\alpha-2)q+2}}{(\alpha-2)q+2}<\infty ,
\]
and the function is integrable on the disk, which is the assertion.

Conversely --- and this is why the range of $q$ in (10) is what it is, though
the statement does not assert it --- for $q\ge2/(2-\alpha)$ one has $s+1\le-1$
and the same computation gives a divergent integral; so the condition
$q<2/(2-\alpha)$ cannot be relaxed.
\end{proof}
